\documentclass[a4paper,11pt]{article}
\usepackage[left=3cm,right=2cm,top=2.7cm,bottom=1.8cm]{geometry}
\usepackage[utf8]{inputenc}
\usepackage[T1]{fontenc}
\usepackage[brazil]{babel}
\usepackage{lmodern,microtype,graphicx,booktabs,tabularx,enumitem,indentfirst}
\usepackage{ragged2e,fancyhdr,titlesec,tikz,hyperref}
\usetikzlibrary{positioning,arrows.meta}
\hypersetup{hidelinks}
\setlength{\parindent}{1.25cm}
\setlength{\headheight}{53pt}
\titleformat{\section}{\normalfont\fontsize{11}{13}\bfseries\MakeUppercase}{\thesection}{0.5em}{}
\titleformat{\subsection}{\normalfont\fontsize{11}{13}\bfseries}{\thesubsection}{0.5em}{}
\titlespacing*{\section}{0pt}{1.2ex}{0.6ex}
\titlespacing*{\subsection}{0pt}{1ex}{0.5ex}
\pagestyle{fancy}
\fancyhf{}
\rhead{\thepage}
\renewcommand{\headrulewidth}{0pt}
\fancypagestyle{firstpage}{
  \fancyhf{}
  \lhead{\includegraphics[width=0.22\textwidth]{Imagens/FURG.png}}
  \chead{\small\textbf{Universidade Federal do Rio Grande -- FURG}\\Sistemas Distribuídos\\Rio Grande/RS, Brasil\\Graduação em Engenharia de Computação}
  \renewcommand{\headrulewidth}{0pt}
}

\begin{document}
\thispagestyle{firstpage}
\justifying
\vspace*{0.5cm}
\begin{center}\textbf{CLIENTE/SERVIDOR COM WEB SERVICE: TRAVELAI}\end{center}
\vspace{0.5cm}
\noindent\textbf{Membros do grupo:} Andrey Vinicius Santos Souza -- 164402; Guilherme Rodrigues Estrella -- 164391.

\section{Motivação e definição do problema}
Ao planejar uma viagem, o usuário costuma consultar sites diferentes para descobrir a localização de uma cidade, a previsão do tempo, dados do país e a cotação da moeda local em relação ao real brasileiro (BRL). Essa consulta fragmentada dificulta a comparação das informações. O TravelAI reúne os dados em uma interface web e trata as respostas de serviços externos antes de apresentá-las.

\section{Arquitetura do sistema e regras de funcionamento}
O navegador executa a interface React e consulta diretamente quatro serviços públicos. A camada \texttt{src/services} valida as respostas JSON com Zod e as converte em modelos internos de \texttt{src/types}, apresentados pelos componentes em vez de expor os JSONs brutos. O roteiro e as perguntas contextuais usam uma função serverless em \texttt{/api/gemini}, que consulta o Google Gemini sem expor a chave de API ao navegador.

\subsection{Tecnologias utilizadas}
React 19, TypeScript, Vite e Tailwind CSS compõem a interface. A Fetch API realiza as requisições; um cliente HTTP comum aplica timeout e converte falhas em erros conhecidos. Vitest e Testing Library cobrem a lógica e os componentes. Playwright testa a navegação no navegador com respostas simuladas. A função serverless usa a API REST do Gemini.

\subsection{Web services integrados}
\begin{table}[ht]
\centering\small
\caption{Serviços e informações usadas no TravelAI}
\begin{tabularx}{\linewidth}{@{}lX@{}}
\toprule
\textbf{Serviço} & \textbf{Informações apresentadas ao usuário} \\
\midrule
Open-Meteo Geocoding & Cidade, país, código ISO, coordenadas e fuso horário. \\
Open-Meteo Forecast & Condições atuais e previsão de cinco dias. \\
countries.dev & Nome, bandeira, capital, população, idiomas e moedas do país. \\
Frankfurter & Cotação indicativa de BRL para a moeda local e data de referência. \\
Google Gemini & Roteiro de viagem e resposta contextual a uma pergunta. \\
\bottomrule
\end{tabularx}
\end{table}

\subsection{Fluxo das requisições}
\begin{figure}[ht]
\centering
\begin{tikzpicture}[node distance=7mm and 7mm, every node/.style={font=\small,align=center}, box/.style={draw,rounded corners,minimum height=8mm,text width=37mm}, arrow/.style={-{Stealth[length=2mm]}}]
\node[box] (ui) {Interface React\\no navegador};
\node[box,right=of ui] (services) {\texttt{src/services}\\validação e normalização};
\node[box,right=of services] (public) {Open-Meteo\\countries.dev\\Frankfurter};
\node[box,below=of services] (proxy) {Função \texttt{/api/gemini}\\chave no servidor};
\node[box,right=of proxy] (gemini) {API Google Gemini};
\draw[arrow] (ui) -- (services);
\draw[arrow] (services) -- (public);
\draw[arrow] (public.north) to[out=125,in=55] node[above,align=center]{JSON} (services.north);
\draw[arrow] (services.north) to[out=125,in=55] node[above,align=center]{modelos internos} (ui.north);
\draw[arrow] (ui.south) |- (proxy.west);
\draw[arrow] (proxy) -- (gemini);
\end{tikzpicture}
\caption{Serviços públicos acessados pelo navegador; Gemini acessado pela função serverless.}
\end{figure}

\newpage
Depois da busca de uma cidade, o usuário seleciona um resultado; entre nomes repetidos, a aplicação prioriza o país esperado e centros urbanos. A seleção fornece coordenadas e fuso para a previsão, e o código ISO do país para consultar os dados nacionais. A cotação parte de BRL e é solicitada apenas quando a moeda do destino é compatível. Clima, país e câmbio têm estados independentes de carregamento, sucesso e falha; uma nova busca cancela as requisições anteriores. O navegador recebe os modelos internos após validação com Zod e apresenta campos formatados, não o JSON bruto.

A Vercel foi escolhida por hospedar os arquivos estáticos do Vite e a função serverless no mesmo domínio HTTPS. Isso oferece acesso remoto pelo navegador e mantém a chave Gemini no ambiente do servidor, sem exigir a administração de um servidor próprio. A interface envia as solicitações de IA como JSON para \texttt{/api/gemini}; a função restringe método e origem, limita o corpo, valida a entrada e solicita JSON estruturado. A resposta é validada novamente e, no roteiro, a quantidade e a sequência dos dias também são conferidas antes de retornar ao frontend. Em falhas transitórias ou timeout, a função tenta um modelo alternativo.

\section{Acesso à aplicação e execução local}
\textbf{Aplicação online:} \url{https://travelai-furg.vercel.app}. A página foi publicada na Vercel e respondeu HTTP 200; a avaliação pode ser feita pelo navegador, sem instalação. \textbf{Código-fonte:} \url{https://github.com/dreyvinixz/TripScope}. Em 25 de outubro de 2026, os fluxos Gemini de pergunta contextual e geração de roteiro responderam HTTP 200. Quando o modelo principal retornou HTTP 503, o fallback \texttt{gemini-3.5-flash-lite} concluiu a solicitação.

Para desenvolvimento local, são necessários Node.js 20 ou superior e npm. Na raiz do repositório, executar \texttt{npm ci} e \texttt{npm run dev}; o Vite informa o endereço local. As quatro APIs públicas não exigem chave. Para usar o Gemini, definir \texttt{GEMINI\_API\_KEY} no ambiente do servidor, com base em \texttt{.env.example}, e reiniciar o Vite. A chave não deve usar o prefixo \texttt{VITE\_}.

Os comandos de verificação são \texttt{npm run lint}, \texttt{npm run typecheck}, \texttt{npm test}, \texttt{npm run build} e \texttt{npm run test:e2e}. Os testes unitários e E2E principais usam \textit{fixtures} e \textit{mocks}. \texttt{npm run test:contract} consulta os serviços públicos e depende da rede; a verificação final deve incluir a aplicação hospedada sem \textit{mocks}.

\section{Dificuldades enfrentadas e soluções adotadas}
\subsection{Indisponibilidade e formatos externos}
Quatro serviços diferentes podem responder lentamente, falhar ou alterar seus formatos. O cliente HTTP converte erros em categorias conhecidas, e Zod verifica os JSONs antes da conversão para os modelos internos. O painel informa a falha do recurso afetado e mantém os demais cartões disponíveis.

\subsection{Proteção da chave do Gemini}
Uma chave inserida no código do cliente ficaria visível no navegador. A integração com Gemini foi isolada em \texttt{/api/gemini}. A função lê a chave do ambiente do servidor, valida a entrada e devolve apenas a resposta tratada para a interface.

\subsection{Testes reproduzíveis}
Testes automatizados que dependem de provedores externos podem falhar por rede ou limite de requisições. A suíte principal usa respostas simuladas e determinísticas; testes de contrato e a demonstração online verificam a integração real separadamente.

\section{Conclusão}
O \textbf{TravelAI} reúne informações de destinos vindas de serviços web distintos em uma interface útil ao viajante. A validação e a normalização das respostas reduzem o acoplamento com os formatos externos, enquanto os estados independentes permitem uso parcial quando um provedor falha. A integração com Gemini mantém a credencial fora do navegador, valida os dados e apresenta respostas estruturadas. Os testes online confirmaram a geração de roteiros e respostas contextuais, incluindo o uso de um modelo alternativo quando o principal ficou indisponível.

\end{document}
